\chapter{Ethical and Legal Considerations}

Intrusion detection tools operate on network traffic that may contain sensitive information, and a system that decides between permitting and blocking carries direct operational consequences. Although this thesis is limited to an experimental prototype for offline inference, the design decisions that make it reproducible and safe also have an ethical and legal dimension that must be made explicit.

\section{Dual Use and Strictly Defensive Scope}
Cybersecurity research has an intrinsically dual nature: the same techniques that allow detecting malicious traffic can, in principle, inform its evasion. The work described here is defensive by construction and, moreover, deliberately inert with respect to any real network. The \texttt{PERMIT} and \texttt{BLOCK} outputs are experimental classifications over static flow records: the system does not perform live blocking, packet dropping, or \texttt{iptables} manipulation, and it does not integrate with firewalls, Software-Defined Networking (SDN) controllers, or Intrusion Prevention Systems (IPS), as delimited in the non-objectives of Chapter 2. This absence of online actuation eliminates the risk of the artifact, by itself, causing undue blocking of legitimate traffic in a production network. The system is an experimental measurement instrument, not a deployable defense mechanism.

\section{Privacy and Data Handling}
The system operates exclusively on flow-level statistical features and not on the payload of the packets, avoiding deep packet inspection and much of its legal and privacy implications. The canonical feature contract reinforces this position by excluding at the schema level potentially sensitive identifiers or re-identifiers---source and destination IP addresses, absolute timestamps, \texttt{Flow ID}, and port numbers---which are discarded before any training. The CICIDS2017 benchmark is a public dataset published by the Canadian Institute for Cybersecurity \parencite{Sharafaldin2018CICIDS2017,CICIDS2017Web}. The traffic used in Phase 2 was generated by the operator in a closed home lab and non-adversarial environment: it is proprietary traffic, contains no personal data from third parties, and was not captured from external networks, so its use poses no issues of consent or exposure of personal information.

\section{Operational Risk and Error Asymmetry}
The consequences of the two types of errors are deeply asymmetric. A false negative---an attack classified as benign and permitted---can enable data exfiltration or service disruption; a false positive---legitimate traffic blocked---interrupts legitimate communications and consumes analysis time. The design of the reward function itself encodes this asymmetry by penalizing false negatives more severely. This thesis reports these risks honestly instead of hiding them behind aggregate accuracy: on the fixed in-distribution random split, the official MAIN model presents a false positive rate of 0.00658 and a false negative rate of 0.00465; however, under the substantially more demanding day split (Check C), the attack recall drops to 0.52954. This performance drop under changing conditions is precisely the reason the system is presented as a research prototype and not as a deployment-ready tool: responsible operational use would require validation on labeled production traffic, multiple capture environments, and an independent security analysis before delegating any decision with real effects to the model \parencite{Apruzzese2022CrossEvaluationNIDS,SommerPaxson2010ClosedWorld}.

\section{Reproducibility as a Responsible Practice}
Finally, the reproducibility discipline described in the methodology is also an ethical practice: security research using machine learning frequently suffers from missing code and irreproducible evaluations, making it difficult to verify or refute published claims \parencite{Olszewski2023ReproSecurityML,Arp2020DosDontsMLSecurity}. The per-run artifact persistence, the fixed test split with seed 42, and its baseline manifest allow any third party to audit exactly under what conditions each figure reported in this thesis was obtained. Making limitations verifiable, rather than presenting them as a statement of good intentions, is the form of methodological honesty that this work pursues.
